\begin{lemma}
For every $\varepsilon>0$ there are integers $\Delta_0,\gamma_0$ such that
the following holds whenever $\Delta\ge\Delta_0$, $\gamma\ge\gamma_0$,
and $0<\gamma\le\Delta$. Let $H$ have finite edge-index set and maximum
degree at most $D$, let $K$ be a finite color type, and let every retained
list $M(e)$ have cardinality $\ell$. Suppose the line graph has degree at
most $\Delta$ and $b_{L(H),\Delta}(e)\le b$ for every edge. For every
completion datum $Q$, use the product experiment with $p=\gamma/\Delta$.
Then
\[
\mathbb EY_x\le D\left(1-\frac{1-e^{-\gamma}}{\Delta}
+\frac2{\Delta^2}\right)^\ell,\qquad
\mathbb E|\operatorname{Del}_e|\le(b+\varepsilon)\ell.
\]
The statistics are those of the selections $I_a$ from this same experiment.
\end{lemma}
\begin{proof}
Take the thresholds in \noderef{IndependentSetSamplingBound} for error
$\varepsilon$. Obtain the completed-set marginal laws and color independence
from \noderef{NibbleProductSamplingLaw}. All the statistics are finite sums
of measurable indicators by \noderef{NibbleResidualDegree} and
\noderef{NibbleDeletedColors}, hence bounded and integrable. By
\noderef{NibbleCompletionData}, each completed graph is $\Delta$-regular
and each $j_a$ is an injective induced embedding. For an original
edge $e$ and $a\in M(e)$, the one-vertex estimate applied at $j_a(e)$ gives
$\Pr[e\in I_a]\ge (1-e^{-\gamma})/\Delta-2/\Delta^2$.
The complement probability is at most
$q=1-(1-e^{-\gamma})/\Delta+2/\Delta^2$.
Here $\Delta\ge1$, so $q\ge1-1/\Delta\ge0$.
Here \noderef{NibbleSelectedEdges} identifies original selection with
completed selection at $j_a(e)$; for $a\notin M(e)$ selection is impossible.
Independence across colors
therefore gives $\Pr[e\text{ unselected}]=\prod_{a\in M(e)}
\Pr[e\notin I_a]\le q^\ell$.
By \noderef{NibbleResidualDegree}, $Y_x$ is the sum of the unselected-edge
indicators over edges through $x$. There are at most $D$ such edge indices
by \noderef{HypergraphDegree} and \noderef{MaxDegreeAtMost}.
Summing their expectations proves
the first inequality, including $\ell=0$.

Fix $e$ and $a\in M(e)$. By \noderef{NibbleDeletedColors} and the pullback
formula in \noderef{NibbleSelectedEdges}, the deletion event for $a$ asks that the completed
selection meet the image $X$ of the neighbors of $e$ whose retained lists
contain $a$. These vertices are neighbors of $j_a(e)$ by the inducedness
in \noderef{NibbleCompletionData}.
The second estimate in \noderef{IndependentSetSamplingBound} bounds its
probability by the local parameter in $Q_a[X\cup\{j_a(e)\}]$ plus
$\varepsilon$. The map $j_a$ bijects this rooted graph with the induced
graph of $L(H)$ on $\{e\}\cup\{f:f\sim e,\ a\in M(f)\}$.
Under this bijection, neighbors correspond, every neighbor in the completed
rooted graph has a preimage, and adjacency is preserved and reflected.
Thus \noderef{LocalBParameterNeighborhoodEmbeddingTransport} gives equality
of the local parameters. Apply \noderef{NibbleInducedBParameterComparison}
to this finite subset of $L(H)$: the parameter is at most
$b_{L(H),\Delta}(e)\le b$. Thus each of the $\ell$ deletion indicators
has mean at most $b+\varepsilon$. Sum them using
\noderef{NibbleDeletedColors}. This proves the second inequality without
any assumption on the sign or size of $b$.
\end{proof}
